\documentclass[aspectratio=169,10pt]{beamer}
% Estilo autónomo: no requiere options.org ni los macros TikZ originales.
\usepackage[T1]{fontenc}
\usepackage[utf8]{inputenc}
\usepackage[spanish]{babel}
\usepackage{lmodern}
\usepackage{amsmath,amssymb}
\usepackage{graphicx}
\usepackage{xcolor}
\makeatletter
\@ifclassloaded{beamer}{%
\usetheme{Boadilla}
\definecolor{AzulTitulo}{HTML}{14253B}
\definecolor{VerdeEstructura}{HTML}{127C84}
\setbeamercolor{frametitle}{fg=white,bg=AzulTitulo}
\setbeamercolor{structure}{fg=VerdeEstructura}
\setbeamertemplate{navigation symbols}{}
\setbeamertemplate{footline}[frame number]
\setbeamersize{text margin left=7mm,text margin right=7mm}
}{}
\makeatother
\newcommand{\AnimWide}[1]{%
  \href{run:videos/#1.mp4?autostart}{%
    \includegraphics[width=.97\linewidth,height=.60\textheight,keepaspectratio]{posters/#1.png}}}
\newcommand{\AnimSide}[1]{%
  \href{run:videos/#1.mp4?autostart}{%
    \includegraphics[width=\linewidth,height=.69\textheight,keepaspectratio]{posters/#1.png}}}
\newcommand{\AnimTitle}[1]{%
  \href{run:videos/#1.mp4?autostart}{%
    \includegraphics[width=.64\linewidth,height=.43\textheight,keepaspectratio]{posters/#1.png}}}
\newcommand{\SmallText}{\fontsize{9.1}{11.3}\selectfont}

\title{Una biblioteca de Python\\para gráficas de clanes y tipo de homotopía}
\author{Rafael Villarroel Flores, UAEH}
\date{22 de octubre de 2026}
\begin{document}
\begin{frame}[plain]
  \centering
  {\Large\bfseries Una biblioteca de Python\\[.45em]para gráficas de clanes y tipo de homotopía\par}
  \medskip
  {Rafael Villarroel Flores, UAEH\par}
  \smallskip
  {22 de octubre de 2026\par}
  \vspace{1em}
  \AnimTitle{intro}
\end{frame}

\section{El operador de clanes en gráficas}

\begin{frame}{Gráficas}
\begin{columns}[T,onlytextwidth]
\column{.53\textwidth}
\vspace{.3em}
Consideramos solamente gráficas finitas y simples (no dirigidas, sin lazos, sin aristas múltiples).
\par\medskip
La cantidad de vértices de la gráfica $G$ se llama el \alert{orden} de $G$ y lo denotamos $|G|$.
\column{.45\textwidth}
\centering\AnimSide{graficas}
\end{columns}
\end{frame}

\begin{frame}{Completas}
\centering
Una \alert{completa} de una gráfica es un conjunto de vértices mutuamente adyacentes.
\par\medskip
\AnimWide{completas}
\end{frame}

\begin{frame}{Clanes}
\centering
Un \alert{clan} de una gráfica es una completa maximal bajo inclusión.
\par\medskip
\AnimWide{clanes}
\end{frame}

\begin{frame}{Gráfica de clanes}
\begin{columns}[T,onlytextwidth]
\column{.57\textwidth}
\SmallText
Dada una gráfica $G$, su \alert{gráfica de clanes $K(G)$} es la gráfica de intersección de sus clanes.
\par\medskip
Identificamos cada clan $q$ con un vértice de $K(G)$:
\begin{itemize}
\item $q\in K(G)$ si y sólo si $q$ es un clan de $G$.
\item Para clanes distintos $q,q'$, hay arista $qq'$ en $K(G)$ si y sólo si $q\cap q'\ne\varnothing$.
\end{itemize}
\column{.41\textwidth}
\centering\AnimSide{grafica_clanes}
\end{columns}
\end{frame}

\begin{frame}{Ejemplo}
\centering
Dos clanes maximales de $G$ se convierten en dos vértices adyacentes de $K(G)$.
\par\medskip
\AnimWide{ejemplo}
\end{frame}

\begin{frame}{Gráficas iteradas de clanes}
\centering
La sucesión de \alert{gráficas iteradas de clanes} se define mediante
\[
K^{0}(G)=G,\qquad K^{n}(G)=K\bigl(K^{n-1}(G)\bigr),\quad n\ge 1.
\]
\vspace{-.4em}
\AnimWide{iteradas}
\end{frame}

\begin{frame}{Prisma}
\centering
El prisma triangular $T_3$ y su gráfica de clanes $K(T_3)$.
\par\medskip
\AnimWide{prisma}
\par\smallskip
En este caso, $K^2(T_3)\cong T_3$.
\end{frame}

\begin{frame}{Suspensión de $C_5$}
\centering
Sea $G$ la suspensión del ciclo de cinco vértices $C_5$.
A la izquierda, $G$; a la derecha, $K(G)$.
\par\medskip
\AnimWide{suspension_c5}
\end{frame}

\begin{frame}{Comportamiento}
\begin{columns}[T,onlytextwidth]
\column{.60\textwidth}
\SmallText
\begin{itemize}
\item $G$ es \alert{convergente} si la sucesión $|K^n(G)|$ es acotada.
\item Equivale a que $G,K(G),K^2(G),\ldots$ tenga sólo una cantidad finita de términos distintos salvo isomorfismo.
\item Equivale a que $K^m(G)\cong K^n(G)$ para algunos $m<n$.
\item Si $G$ no es convergente, es \alert{divergente}.
\item Decidir el comportamiento es un problema central del área.
\end{itemize}
\column{.38\textwidth}
\centering\AnimSide{comportamiento}
\end{columns}
\end{frame}

\section{Teoremas sobre comportamiento}

\begin{frame}{Iteradas del octaedro}
\centering
Si definimos el \alert{octaedro} $O_n=\overline{nK_2}$, entonces
\[
K(O_n)\cong O_{2^{n-1}}.
\]
Por tanto $O_n$ es divergente si $n\ge3$. (Neumann-Lara, 1973.)
\par\medskip
\AnimWide{iteradas_octaedro}
\end{frame}

\begin{frame}{Gráficas Helly}
\begin{columns}[T,onlytextwidth]
\column{.60\textwidth}
\SmallText
Una colección $\mathcal{C}$ de subconjuntos de $X$ es \alert{intersecante} si dos cualesquiera de sus miembros tienen intersección no vacía.
\par\medskip
Es \alert{Helly} si toda subfamilia intersecante tiene intersección total no vacía.
\par\medskip
Una gráfica $G$ es \alert{Helly} si su colección de clanes satisface la propiedad de Helly.
\column{.38\textwidth}
\centering\AnimSide{helly}
\end{columns}
\end{frame}

\begin{frame}{La pirámide y propiedad de Helly}
\begin{columns}[T,onlytextwidth]
\column{.58\textwidth}
\SmallText
\begin{itemize}
\item La \alert{pirámide} es la gráfica no Helly más pequeña.
\item Las gráficas sin triángulos son Helly. También lo son $K_4$ y $T_3$.
\item Existe un algoritmo polinomial que decide si una gráfica es Helly (Dragan, 1989; Szwarcfiter, 1997).
\end{itemize}
En la animación se distinguen tres clanes que se intersectan dos a dos, pero no comparten un vértice común.
\column{.40\textwidth}
\centering\AnimSide{piramide}
\end{columns}
\end{frame}

\begin{frame}{Teoremas de Escalante}
\begin{columns}[T,onlytextwidth]
\column{.60\textwidth}
\SmallText
\begin{itemize}
\item (Escalante, 1973) Si $G$ es Helly, entonces $G$ es convergente.
\item (Escalante, 1973) Si $G$ es Helly, $K(G)$ es Helly.
\item Existen gráficas no Helly cuyo gráfico de clanes sí es Helly, como la pirámide.
\item De aquí se obtiene que toda gráfica de a lo sumo cinco vértices es convergente.
\end{itemize}
\column{.38\textwidth}
\centering\AnimSide{escalante}
\end{columns}
\end{frame}

\begin{frame}{Gráficas desmantelables}
\begin{columns}[T,onlytextwidth]
\column{.60\textwidth}
\SmallText
Un vértice $x$ de $G$ es \alert{dominado} si existe $y$ adyacente a $x$ tal que todo vecino de $x$ distinto de $y$ es adyacente a $y$.
\par\medskip
La gráfica $G$ es \alert{desmantelable} si puede reducirse a la gráfica de un solo vértice mediante remociones sucesivas de vértices dominados.
\par\medskip
Los árboles y las gráficas cordales conexas son desmantelables.
\column{.38\textwidth}
\centering\AnimSide{desmantelables}
\end{columns}
\end{frame}

\begin{frame}{Teoremas sobre vértices dominados}
\begin{columns}[T,onlytextwidth]
\column{.61\textwidth}
\SmallText
\begin{itemize}
\item (Prisner, 1992) Si $G$ es desmantelable, existe $n$ tal que $K^n(G)$ es la gráfica de un vértice: $G$ es \alert{nula}.
\item (Frías-Armenta, Neumann-Lara, Pizaña, 2004) Si $x$ es dominado en $G$, entonces $G$ y $G-x$ tienen el mismo comportamiento.
\item Por ejemplo, el octaedro $O_3$ con dos hojas adjuntas sigue siendo divergente.
\end{itemize}
\column{.37\textwidth}
\centering\AnimSide{dominados}
\end{columns}
\end{frame}

\end{document}
